\subsection{Kan complexes}

Recall that $\sSet=[\Delta^\op,\Set]$.
Since $\Set$ is complete and cocomplete, so too is $\sSet$.
That is to say, $\sSet$ has all (small) limits and colimits.
The co/limit of a functor is taken coordinate-wise, so such limits and colimits are easy to describe.

Most important for now is the product of simplicial sets.
If $K,L$ are simplicial sets, then their product $K\times L$ is the simplicial set where
$$ (K\times L)_n = K_n\times L_n\hbox{ for $n<\omega$},\qquad (K\times L)\phi = (K\phi)\times(L\phi)\hbox{ for $\phi\colon[n]\to[m]$} . $$
Explicitly, the face and degeneracy maps are
$$ d_i(x,y) = (d_i(x),d_i(y)),\qquad s_i(x,y) = (s_i(x),s_i(y)) . $$

If $X$ is a simplicial set, a {\emphcolor simplicial subset} of $X$ is a simplicial set $Y$ such that $Y_n\subseteq X_n$ for all $n$ and the face and degeneracy maps
of $X$ restrict to those of $Y$.
This is equivalent to $Y$ being a subobject of $X$ in the functor category $\sSet=[\Delta^\op,\Set]$ (see the following note).

\bnote

    In a functor category $[\C,\D]$ there are two notions of a ``subfunctor''.
    Let $F,G\colon\C\to\D$ be functors, then we may say that $F$ is a subfunctor of $G$ iff it is a subobject of $G$ in the functor category.
    This means that there is a monic natural transformation $F\To G$.

    Another, generally easier to validate, definition is that $F$ is a subfunctor of $G$ iff $F(c)$ is a subobject of $G(c)$ in a natural way.
    That is to say, there is a natural transformation $\alpha\colon F\To G$ whose components $\alpha_c\colon F(c)\to G(c)$ are monic.

    In general clearly the latter implies the former: a pointwise monic natural transformation is clearly monic.
    Under relatively weak conditions, that $\D$ has pullbacks, the former also implies the latter.

    Recall that a morphism $f\colon X\to Y$ is monic iff its pullback with itself is $X$ (with projections given by the identity).
    That is, iff the following square is a pullback

    \commsq{X&X\cr X&Y}
        {\id}{f}{\id}{f}

    Thus if $\alpha\colon F\To G$ is monic, then its pullback with itself is $F$.
    Limits in the functor category are computed pointwise, so this means that $\alpha_c$'s pullback with itself is $Fc$.
    This means that $\alpha_c$ is monic for all $c$, as desired.

\enote

Notice that if $\set{X^i}_{i\in I}$ is a collection of simplicial subsets of $X$, then so is their intersection and union.
This is because the $X^i$ are invariant under any $\phi\colon X_m\to X_n$, and thus so too is their intersection and union.

\bdefn

    If $S\subseteq X$ is an arbitrary subset (meaning a subset of $\bigcup_nX_n$), then the simplicial subset it {\emphcolor generates} is the smallest
    simplicial subset of $X$ containing $S$ (meaning it contains $S\cap X_n$ as $n$-simplicies).
    Such a minimal simplicial subset exists as we can take the intersection of all simplicial subsets containing $S$.

\edefn

For a morphism in the simplex category $\phi\colon[n]\to[m]$, let us write $X_\phi\colon X_m\to X_n$ for the induced map.
Then if $S\subseteq X$ is an arbitrary set, then the simplicial subset it generates is just
$$ (S) = \set{X_\phi(x)}[x\in S] . $$
This is a not-so-formal way of saying that $(S)_n$ is the collection of all $X_\phi(x)$ where $x\in S\cap X_k$ and $\phi\colon[n]\to[k]$.

Recall the standard $n$-simplex $\Delta^n=\hom_\Delta(\bul,[n])$.
Its $k$-simplicies $\Delta^n_k$ are then order-preserving maps $[k]\to[n]$.
In particular, $\Delta^n_{n-1}$ consists only of $\delta^0,\dots,\delta^n$.
Recall that the $k$-simplex in $\Delta^n$ corresponding to $\phi\colon[k]\to[n]$ is given by $[\phi(0),\dots,\phi(k)]$.
In particular $\delta^i$ corresponds to $[0,\dots,\hat k,\dots,n]$.

\bdefn

    The {\emphcolor boundary of $\Delta^n$}, $\partial\Delta^n$, is the simplicial subset of $\Delta^n$ generated by $\delta^0,\dots,\delta^n$.

    The {\emphcolor simplicial $n$-horn} $\Lambda^n_k$ is the simplicial subset of $\Delta^n$ generated by $\delta^0,\dots,\hat\delta^k,\dots,\delta^n$
    (all coface maps except for $\delta^k$).

\edefn

The $i$th face of $\Delta^n$ is the simplicial subset $\partial_i\Delta^n\subseteq\Delta^n$ generated by $\delta^i$.
Since the union of simplicial subsets is a simplicial subset, we see that
$$ \partial\Delta^n = \bigcup_{i\leq n}\partial_i\Delta^n,\qquad \Lambda^n_k = \bigcup_{i\neq k}\partial_i\Delta^n . $$
In more generality, for a set of indices $I\subseteq[n]$ we can define $\Delta^I$ to be the simplicial subset of $\Delta^n$ consisting of functions $f\colon[k]\to[n]$ whose
image lies in $I$.
This is indeed a simplicial subset because precomposing $f$ with $\phi\colon[m]\to[k]$ still has its image lie in $I$.

Then $\partial_i\Delta^n=\Delta^{[n]-i}$.
Indeed, $\delta^i$ by definition ``misses'' $i$.
And if $\phi\colon[k]\to[n]$ misses $i$, then we can find $\psi\colon[k]\to[n-1]$ such that $\phi=\delta^i\circ\psi$ so that $\phi$ is generated by $\delta^i$.
Then
$$ \partial\Delta^n = \bigcup_{i=0}^n\Delta^{[n]-i} $$
so that its $k$th simplices consists of maps $\phi\colon[k]\to[n]$ which are not surjective (as it must miss some $i$).

Similarly
$$ \Lambda^n_k = \bigcup_{i\neq k}\Delta^{[n]-i},\qquad (\Lambda^n_k)_j = \set{f\colon[j]\to[n]}[\hbox{$f$ does not have $0,\dots,\hat k,\dots,n$ in its image}] . $$
Since $\delta^k$ is the unique map without $k$ in its image, geometrically we view $\Lambda^n_k$ as the union of the faces of $\Delta^n$
except for the one opposite the $k$th vertex.
For $n=2$ for example, the three horns are

\bigskip
\hbox to\hsize{\hfil%
    \def\diagrowheight{20pt}\def\diagcolwidth{10pt}%
    \vbox{\halign{\hfil#\hfil\cr\drawdiagram{&$1$\cr$0$&&$2$}{\diagarrow{from={2,1}, to={2,3}}\diagarrow{from={2,1}, to={1,2}}}\cr$\Lambda^2_0$\cr}}\hfil%
    \vbox{\halign{\hfil#\hfil\cr\drawdiagram{&$1$\cr$0$&&$2$}{\diagarrow{from={1,2}, to={2,3}}\diagarrow{from={2,1}, to={1,2}}}\cr$\Lambda^2_1$\cr}}\hfil%
    \vbox{\halign{\hfil#\hfil\cr\drawdiagram{&$1$\cr$0$&&$2$}{\diagarrow{from={2,1}, to={2,3}}\diagarrow{from={1,2}, to={2,3}}}\cr$\Lambda^2_2$\cr}}\hfil%
}

Note that there exists a coequalizer diagram
$$ \coprod_{0\leq i<j\leq n}\Delta^{n-2} \tto \coprod_{i=0}^n\Delta^{n-1} \to \partial\Delta^n $$
where the first morphisms have coordinates $\Delta^{n-2}\to\coprod_i\Delta^{n-1}$ correspond to $\delta^{j-1}$ and $\delta^i$
respectively.
That is,

\bigskip
\centerline{\drawdiagram{
    $\Delta^{n-2}$&$\Delta^{n-1}$\cr
    $\displaystyle\coprod_{i<j}\Delta^{n-2}$&$\displaystyle\coprod_i\Delta^{n-1}$&$\partial\Delta^n$\cr
    $\Delta^{n-2}$&$\Delta^{n-1}$\cr
}{
    \diagarrow{from={1,1}, to={1,2}, text={\delta^{j-1}}, y distance=.25cm}
    \diagarrow{from={1,1}, to={2,1}, text={\inc_{i<j}}, x distance=-.25cm}
    \diagarrow{from={3,1}, to={2,1}, text={\inc_{i<j}}, x distance=-.25cm}
    \diagarrow{from={3,1}, to={3,2}, text={\delta^i}, y distance=-.25cm}
    \diagarrow{from={1,2}, to={2,2}, text={\inc_i}, x distance=.25cm}
    \diagarrow{from={3,2}, to={2,2}, text={\inc_j}, x distance=.25cm}
    \diagarrow{from={1,2}, to={2,3}, text={\delta^i}, x distance=.25cm}
    \diagarrow{from={3,2}, to={2,3}, text={\delta^j}, x distance=.25cm}
    \diagarrow{from={2,1}, to={2,2}, y off=2pt}
    \diagarrow{from={2,1}, to={2,2}, y off=-2pt}
    \diagarrow{from={2,2}, to={2,3}}
}}

Similarly there is a coequalizer

\bigskip
\centerline{\drawdiagram{
    $\Delta^{n-2}$&$\Delta^{n-1}$\cr
    $\displaystyle\coprod_{i<j}\Delta^{n-2}$&$\displaystyle\coprod_{i\neq k}\Delta^{n-1}$&$\Lambda^n_k$\cr
    $\Delta^{n-2}$&$\Delta^{n-1}$\cr
}{
    \diagarrow{from={1,1}, to={1,2}, text={\delta^{j-1}}, y distance=.25cm}
    \diagarrow{from={1,1}, to={2,1}, text={\inc_{i<j}}, x distance=-.25cm}
    \diagarrow{from={3,1}, to={2,1}, text={\inc_{i<j}}, x distance=-.25cm}
    \diagarrow{from={3,1}, to={3,2}, text={\delta^i}, y distance=-.25cm}
    \diagarrow{from={1,2}, to={2,2}, text={\inc_i}, x distance=.25cm}
    \diagarrow{from={3,2}, to={2,2}, text={\inc_j}, x distance=.25cm}
    \diagarrow{from={1,2}, to={2,3}, text={\delta^i}, x distance=.25cm}
    \diagarrow{from={3,2}, to={2,3}, text={\delta^j}, x distance=.25cm}
    \diagarrow{from={2,1}, to={2,2}, y off=2pt}
    \diagarrow{from={2,1}, to={2,2}, y off=-2pt}
    \diagarrow{from={2,2}, to={2,3}}
}}

Using the universal property of coequalizers, the following lemma is immediate.

\blemm

    Let $X$ be an arbitrary simplicial set.
    \benum
        \item A simplicial map $\partial\Delta^n\to X$ is the same as a tuple $(y_0,\dots,y_n)$ of $(n-1)$-simplices
        of $X$ such that $d_iy_j=d_{j-1}y_i$ for $i<j$.
        \item A simplicial map $\Lambda^n_k\to X$ is the same as a tuple $(y_0,\dots,\hat y_k,\dots,y_n)$ of $(n-1)$-simplices
        of $X$ such that $d_iy_j=d_{j-1}y_i$ for $i<j$, neither equal to $k$.
    \eenum

\elemm

\bproof

    A map out of a coequalizer $A\tto B$ is just a map $B\to X$ whose composition with either arrow in the coequalizer
    diagram is equal.
    Couple this with the fact that $\hom_\sSet(\Delta^n,X)\cong X_n$ and you have the desired result.
    \qed

\eproof

\bdefn

    Say that a simplicial set $X$ has dimension $\leq n$ if all $k$-simplices are degenerate for $k>n$.
    The {\emphcolor $n$-skeleton} of $X$ is the largest simplicial subset of $X$ with dimension $\leq n$ (which exists as we can take unions).
    We denote it by $\sk_n(X)$.

\edefn

Alternatively, the $n$-skeleton of $X$ is the simplicial subset generated by all non-degenerate simplices of dimension $\leq n$.

For example, $\sk_{n-1}(\Delta^n)=\partial\Delta^n$.
This is because $\partial\Delta^n_k$ consists only of degenerate simplices for $k\geq n$, and for $k<n$ we have $\partial\Delta^n_k=\Delta^n_k$.

In general we have a filtration of a simplicial set through its skeletons,
$$ \emptyset = \sk_{-1}(X) \subseteq \sk_0(X) \subseteq \sk_1(X) \subseteq \cdots,\qquad X = \bigcup_n\sk_n(X) . $$

\bdefn

    A simplicial map $p\colon X\to Y$ is a {\emphcolor fibration} if for every commutative diagram of solid arrows,
    there is a dashed arrow making it commute

    \centerline{\drawdiagram{
        $\Lambda^n_k$&$X$\cr
        $\Delta^n$&$Y$
    }{
        \diagarrow{from={1,1}, to={1,2}}
        \diagarrow{from={1,2}, to={2,2}}
        \diagarrow{from={1,1}, to={2,1}, left cap={(}}
        \diagarrow{from={1,2}, to={2,2}, text={p}, x distance=.25cm}
        \diagarrow{from={2,1}, to={2,1}, to={2,2}}
        \diagarrow{from={2,1}, to={1,2}, dashed}
    }}

    The arrow $\Lambda^n_k\inj\Delta^n$ is the inclusion map.

    Fibrations are also called {\emphcolor Kan fibrations}.

\edefn

Equivalently, for every tuple of $(n-1)$-simplices $(x_0,\dots,\hat x_k,\dots,x_n)$ such that $d_ix_j=d_{j-1}x_i$ for $i<j$
not equal to $k$, if there is an $n$-simplex $y\in Y_n$ such that $d_iy=p(x_i)$, then there is an $n$-simplex $x\in X_n$
such that $d_ix=x_i$ and $p(x)=y$.

There is a similar concept of fibration for topological spaces.
A continuous map $f\colon T\to U$ is a {\emphcolor Serre fibration} if a dashed arrow exists for any commutative diagram
of continuous maps

\centerline{\drawdiagram{
    $\abs{\Lambda^n_k}$&$T$\cr
    $\abs{\Delta^n}$&$U$
}{
    \diagarrow{from={1,1}, to={1,2}}
    \diagarrow{from={1,2}, to={2,2}}
    \diagarrow{from={1,1}, to={2,1}, left cap={(}}
    \diagarrow{from={1,2}, to={2,2}, text={f}, x distance=.25cm}
    \diagarrow{from={2,1}, to={2,1}, to={2,2}}
    \diagarrow{from={2,1}, to={1,2}, dashed}
}}

Because $\S(\abs{\Delta^n})=\Delta^n$ and $\S(\abs{\Lambda^n_k})=\Lambda^n_k$, and geometric realization is left-adjoint to
$\S$, such diagrams can be identified with diagrams

\centerline{\drawdiagram{
    $\Lambda^n_k$&$\S(T)$\cr
    $\Delta^n$&$\S(U)$
}{
    \diagarrow{from={1,1}, to={1,2}}
    \diagarrow{from={1,2}, to={2,2}}
    \diagarrow{from={1,1}, to={2,1}, left cap={(}}
    \diagarrow{from={1,2}, to={2,2}, text={\S(f)}, x distance=.25cm}
    \diagarrow{from={2,1}, to={2,1}, to={2,2}}
    \diagarrow{from={2,1}, to={1,2}, dashed}
}}

So that $f\colon T\to U$ is a Serre fibration iff $\S(f)\colon\S(T)\to\S(U)$ is a (Kan) fibration.

Notice that for any topological space $X$ and one-point space $*$, $X\to*$ is a Serre fibration.
This is because $\abs{\Lambda^n_k}\inj\abs{\Delta^n}$ is a strong deformation retract.
Thus $\S(X)\to\S(*)$ is a fibration of simplicial sets.
Note that $\S(*)=\Delta^0=*$ is the terminal simplicial set (with a single simplex of each dimension).

This leads to the following definition.

\bdefn

    A simplicial set $X$ is {\emphcolor fibrant}, or a {\emphcolor Kan complex}, if the map $X\to*$ is fibrant.

\edefn

Simplicial sets which arise from geometric simplicial complexes are not Kan complexes, thus $\Delta^n$ is not.
Neither is $\partial\Delta^n$.

The following lemma characterizes equivalent conditions for a simplicial set being fibrant, and is immediate.

\blemm[name={kancompequiv}]

    Let $X$ be a simplicial set, then the following are equivalent.
    \benum
        \item $X$ is fibrant.
        \item Every map $f\colon\Lambda^n_k\to X$ extends to a map $g\colon\Delta^n\to X$ for all $n,k$.

        In a commutative diagram:

        \centerline{\drawdiagram{
            $\Lambda^n_k$&$X$\cr
            $\Delta^n$
        }{
            \diagarrow{from={1,1}, to={1,2}}
            \diagarrow{from={1,1}, to={2,1}, left cap={(}}
            \diagarrow{from={2,1}, to={1,2}, dashed}
        }}

        \item For every tuple of $(n-1)$-simplices $(x_0,\dots,\hat x_k,\dots,x_n)$ of $X$ such that $d_ix_j=d_{j-1}x_i$
        for $i<k$ not $k$, there is an $n$-simplex $x\in X_n$ such that $d_ix=x_i$.
    \eenum

\elemm

Fibrant simplicial sets are extremely important, as we will see soon.

An interesting result is that the underlying simplicial set of every simplicial group is fibrant.
A simplicial group is a simplicial object in the category of groups, i.e.\ a functor $\Delta^\op\to\Grp$.
By underlying simplicial set, we mean the composition of the functor with the forgetful functor $\Grp\to\Set$.

\blemm[title={Moore}]

    If $H\colon\Delta^\op\to\Grp$ is a simplicial group, its underlying set is fibrant.

\elemm

